\documentclass[11pt]{article}
\usepackage[a4paper,margin=0.78in]{geometry}
\usepackage{amsmath,amssymb,amsthm,mathtools}
\usepackage{graphicx,microtype,enumitem,xcolor,hyperref}
\hypersetup{colorlinks=true,linkcolor=blue!45!black,urlcolor=blue!55!black,citecolor=blue!45!black}
\setlist{nosep,leftmargin=1.7em}
\newtheorem{theorem}{Theorem}
\newtheorem{lemma}{Lemma}
\newcommand{\B}{\mathcal B}
\newcommand{\K}{\mathcal K}
\newcommand{\D}{\mathbb D}
\newcommand{\T}{\mathbb T}
\newcommand{\ellp}{\ell_p}
\newcommand{\ellq}{\ell_q}
\definecolor{answer}{RGB}{35,80,125}

\title{SOT$^*$-density of $\ell_p$ contractions with compact-free commutant\\
\large A substantial partial answer to Question 6.4 of arXiv:2105.04635}
\author{Automated mathematical discovery run \texttt{fa\_banach\_001}, lane 0}
\date{26 August 2026}

\begin{document}
\maketitle

\begin{center}
\fcolorbox{answer}{answer!6}{\parbox{0.91\linewidth}{
\textbf{Candidate substantial partial result.}
For every complex $\ell_p$, $1<p<\infty$, contractions which commute with no
nonzero compact operator are SOT$^*$-dense in the unit ball of
$\B(\ell_p)$.  This proves density, not the comeagerness requested by the
word ``typical'' in the source question.  The stronger Lomonosov-hypothesis
clause is not addressed.
}}
\end{center}

\section*{Source question}
Grivaux and Matheron ask in Question 6.4 of \cite{GM}:
\begin{quote}
Let $X=\ell_p$, $1<p<\infty$. Is it true that a typical
$T\in\B_1(X)$ does not satisfy the assumptions of the Lomonosov Theorem?
Is it at least true that a typical $T\in\B_1(X)$ does not commute with a
non-zero compact operator?
\end{quote}
Here ``typical'' means comeager in the SOT$^*$ topology on the contraction
ball.  The source recalls that the fallback assertion is known for $p=2$
from \cite[Theorem 7.5]{GMM}.

\section*{Definitions}
For a reflexive Banach space $X$, SOT$^*$ convergence $T_j\to T$ means
$T_jx\to Tx$ for every $x\in X$ and $T_j^*x^*\to T^*x^*$ for every
$x^*\in X^*$.  We write
\[
 \{T\}'=\{R\in\B(X):RT=TR\}
\]
for the commutant.  The coordinate projections onto the first $N$
coordinates of $\ell_p$ and $\ell_q$ are both denoted by $P_N$.

\section*{Proof intuition}
An SOT$^*$ neighborhood only monitors finitely many vectors and
functionals.  Thus, after a sufficiently large finite coordinate block, one
may install a unilateral forward-shift tail.  Couple the finite block into
the first tail coordinate through a tiny cyclic functional.  The coupling is
small enough to preserve contractivity, yet it makes the eigenspaces of the
adjoint one-dimensional and their eigenvectors total.

Any operator commuting with this adjoint must act on those eigenspaces by a
bounded holomorphic scalar $\varphi$.  As the eigenvalue approaches any
point of the unit circle, the geometric tail of the eigenvector blows up;
after normalization, the eigenvector tends weakly to zero.  A compact
commutant sends it to zero in norm, forcing every radial boundary value of
$\varphi$ to be zero.  Bounded analytic uniqueness gives $\varphi=0$, and
totality then gives that the compact commutant itself is zero.

\clearpage
\section*{Original source image}
\begin{center}
\includegraphics[width=0.99\textwidth,height=0.78\textheight,keepaspectratio]{figures/open_problem_crop.png}
\end{center}
\noindent\footnotesize Source: Grivaux--Matheron, arXiv:2105.04635,
PDF p.~18, Question 6.4.  This is a direct full-width crop from
\texttt{source\_paper.pdf}.
\normalsize

\section*{The density theorem}
\begin{theorem}\label{thm:main}
Let $1<p<\infty$ and let $X=\ell_p$ over the complex scalars.  Then
\[
 \bigl\{T\in\B_1(X):\{T\}'\cap\K(X)=\{0\}\bigr\}
\]
is SOT$^*$-dense in $\B_1(X)$.
\end{theorem}

\begin{proof}
Let $T_0\in\B_1(X)$ and fix a basic SOT$^*$ neighborhood of $T_0$.
Because $P_N\to I$ strongly on both $\ell_p$ and its dual $\ell_q$, the
finite compressions $P_NT_0P_N$ converge to $T_0$ in SOT$^*$.  Choose $N$
large.  On the finite-dimensional space $E=P_NX$, first scale the
compression slightly and then perturb it slightly to an operator
$A\in\B(E)$ such that
\[
 r:=\|A\|<1
 \quad\hbox{and}\quad
 A^*\text{ has pairwise distinct eigenvalues}.
\]
Matrices with simple spectrum are dense, so this can be done while retaining
membership in the prescribed neighborhood after the construction below.

Choose a norm-one cyclic vector $v^*\in E^*$ for $A^*$.  Such a vector exists:
if $u_1^*,\dots,u_d^*$ are eigenvectors for the distinct eigenvalues, any
linear combination having all coefficients nonzero is cyclic.  Decompose
$X=E\oplus_p F$, identify $F$ with $\ell_p(\mathbb Z_+)$, and let
\[
 S(y_0,y_1,\dots)=(0,y_0,y_1,\dots)
\]
be the forward shift on $F$.  Choose $\delta>0$ so small that
\[
 r^p+\delta^p\leq1
\]
and so that the coupling below is invisible to the chosen neighborhood.
Define
\begin{equation}\label{eq:T}
 T(x,y)=\bigl(Ax,\ \delta v^*(x)e_0+Sy\bigr),
 \qquad x\in E,\ y\in F.
\end{equation}
The three displayed output blocks have disjoint coordinate supports, hence
\[
 \|T(x,y)\|_p^p
 =\|Ax\|_p^p+\delta^p|v^*(x)|^p+\|y\|_p^p
 \leq(r^p+\delta^p)\|x\|_p^p+\|y\|_p^p.
\]
Thus $T$ is a contraction.  Moreover, $T$ can be made arbitrarily close to
$T_0$ in SOT$^*$: on the prescribed primal test vectors the shift sees only
their small tails and the new coupling has norm at most $\delta$; on the dual
test vectors the same assertion follows from their small $\ell_q$ tails.
This proves the required approximation once we show that $T$ has no nonzero
compact operator in its commutant.

Let $q=p/(p-1)$ and let $B$ be the backward shift on $\ell_q$.  From
\eqref{eq:T},
\[
 T^*(u,z)=\bigl(A^*u+\delta z_0v^*,\ Bz\bigr),
 \qquad u\in E^*,\ z\in\ell_q.
\]
For $\lambda\in\Omega:=\D\setminus\sigma(A^*)$, put
\[
 k_\lambda=(1,\lambda,\lambda^2,\dots)\in\ell_q,
 \qquad
 h_\lambda=
 \bigl(\delta(\lambda I-A^*)^{-1}v^*,\ k_\lambda\bigr).
\]
Then $T^*h_\lambda=\lambda h_\lambda$.  Conversely,
$Bz=\lambda z$ forces $z$ to be a scalar multiple of $k_\lambda$, and the
first coordinate equation then determines $u$.  Therefore
\begin{equation}\label{eq:onedim}
 \ker(T^*-\lambda I)=\mathbb C h_\lambda
 \qquad(\lambda\in\Omega).
\end{equation}

The family $(h_\lambda)_{\lambda\in\Omega}$ is total in $X^*$.  Indeed,
suppose $(a,y)\in E\oplus_pF$ is annihilated by every $h_\lambda$.  Then
\begin{equation}\label{eq:meromorphic}
 \delta\bigl[(\lambda I-A^*)^{-1}v^*\bigr](a)
 +\sum_{n\geq0}y_n\lambda^n=0
 \qquad(\lambda\in\Omega).
\end{equation}
Let $Q_j$ be the spectral projections of $A^*$ associated with its distinct
eigenvalues $\mu_j$.  Comparing residues at $\mu_j$ in
\eqref{eq:meromorphic} gives $(Q_jv^*)(a)=0$ for every $j$.  Cyclicity says
that the nonzero vectors $Q_jv^*$ form a basis of $E^*$, so $a=0$.  The
remaining power series is identically zero, whence $y=0$.

Now let $K\in\K(X)$ commute with $T$.  By Schauder's theorem, $K^*$ is
compact, and it commutes with $T^*$.  Equation \eqref{eq:onedim} yields
\[
 K^*h_\lambda=\varphi(\lambda)h_\lambda
 \qquad(\lambda\in\Omega).
\]
The scalar function $\varphi$ is holomorphic on $\Omega$: it is obtained by
evaluating the first tail coordinate of $K^*h_\lambda$, since that coordinate
of $h_\lambda$ is $1$.  Also
$|\varphi(\lambda)|\leq\|K^*\|$, because $\varphi(\lambda)$ is an eigenvalue
of $K^*$.  Its finitely many singularities are therefore removable, so
$\varphi\in H^\infty(\D)$.

Fix $\zeta\in\T$ and let $\lambda=r\zeta$ tend radially to $\zeta$.  The
finite-dimensional component of $h_{r\zeta}$ stays uniformly bounded, since
$\|A\|<1$, whereas
\[
 \|k_{r\zeta}\|_q=(1-r^q)^{-1/q}\longrightarrow\infty.
\]
Consequently $u_r:=h_{r\zeta}/\|h_{r\zeta}\|$ is bounded and converges
coordinatewise to zero.  Since finitely supported vectors are dense in
$(\ell_q)^*=\ell_p$, it follows that $u_r\rightharpoonup0$ in $X^*$.
Compactness gives $\|K^*u_r\|\to0$, while the eigenvector identity gives
\[
 \|K^*u_r\|=|\varphi(r\zeta)|.
\]
Thus $\varphi$ has radial limit zero at every $\zeta\in\T$.  Since it is
bounded, dominated convergence in the Cauchy coefficient formula forces
every Taylor coefficient of $\varphi$ to vanish.  Hence $\varphi=0$.
The operator $K^*$ therefore annihilates the total family
$\{h_\lambda:\lambda\in\Omega\}$, so $K^*=0$ and $K=0$.
\end{proof}

\section*{Verification and exact remaining gap}
The proof uses four independent checks:
\begin{enumerate}
\item contractivity is the displayed disjoint-support $p$-norm identity;
\item SOT$^*$ density uses strong convergence of coordinate projections on
both $\ell_p$ and $\ell_q$;
\item totality follows from residues at the simple finite-block eigenvalues
and then uniqueness of a scalar power series;
\item compactness turns radial weak convergence into norm convergence, while
bounded analytic uniqueness converts the boundary vanishing into
$\varphi=0$.
\end{enumerate}

\paragraph{Why this is not yet the full source answer.}
The theorem proves density, whereas ``typical'' requires a comeager set.
The Hilbert proof in \cite{GMM} inserts a dense class into a specially
constructed dense $G_\delta$ by showing
$\|R\|=\|R\|_{\mathrm e}$ for every commutant element $R$; its key model
estimate uses the Hilbert-space von Neumann inequality.  The boundary
argument above proves that compact commutant elements vanish, but it does not
prove that norm--essential-norm equality for arbitrary noncompact commutant
elements on $\ell_p$.  Density by itself cannot be promoted to
comeagerness.  The first, stronger clause concerning the full Lomonosov
hypothesis is also untouched.

\paragraph{Novelty audit.}
The run indexes were searched for arXiv:2105.04635, arXiv:2012.02016,
``typical contraction'', ``compact operator'', ``commutant'', and
``Lomonosov''.  Bounded primary-source arXiv/web searches through 26 August
2026 used the exact question phrase and close SOT$^*$/$\ell_p$ variants.
They recovered the source paper, the earlier Hilbert-space theorem
\cite{GMM}, and later papers on positive contractions, but no theorem giving
this all-$p$ density statement or resolving Question 6.4 for $p\ne2$.
Novelty confidence is moderate, not definitive.

\paragraph{Human review recommendation.}
Immediate expert review is recommended.  Priority checks are the residue
argument establishing totality, the radial weak-null assertion in $\ell_q$,
and the distinction between this density theorem and the still-missing Baire
upgrade.

\begin{thebibliography}{9}
\bibitem{GM} S. Grivaux and É. Matheron,
\emph{Local spectral properties of typical contractions on $\ell_p$-spaces},
arXiv:2105.04635, especially Question 6.4 on PDF p.~18.
\bibitem{GMM} S. Grivaux, É. Matheron, and Q. Menet,
\emph{Does a typical $\ell_p$-space contraction have a non-trivial invariant
subspace?}, arXiv:2012.02016, especially Theorem 7.5 and its proof.
\end{thebibliography}

\end{document}
